% Options for packages loaded elsewhere
\PassOptionsToPackage{unicode}{hyperref}
\PassOptionsToPackage{hyphens}{url}
\PassOptionsToPackage{dvipsnames,svgnames,x11names}{xcolor}
\documentclass[
  10pt,
  fontset=fandol]{ctexart}
\usepackage{xcolor}
\usepackage[margin=16mm]{geometry}
\usepackage{amsmath,amssymb}
\setcounter{secnumdepth}{-\maxdimen} % remove section numbering
\usepackage{iftex}
\ifPDFTeX
  \usepackage[T1]{fontenc}
  \usepackage[utf8]{inputenc}
  \usepackage{textcomp} % provide euro and other symbols
\else % if luatex or xetex
  \usepackage{unicode-math} % this also loads fontspec
  \defaultfontfeatures{Scale=MatchLowercase}
  \defaultfontfeatures[\rmfamily]{Ligatures=TeX,Scale=1}
\fi
\usepackage{lmodern}
\ifPDFTeX\else
  % xetex/luatex font selection
\fi
% Use upquote if available, for straight quotes in verbatim environments
\IfFileExists{upquote.sty}{\usepackage{upquote}}{}
\IfFileExists{microtype.sty}{% use microtype if available
  \usepackage[]{microtype}
  \UseMicrotypeSet[protrusion]{basicmath} % disable protrusion for tt fonts
}{}
\makeatletter
\@ifundefined{KOMAClassName}{% if non-KOMA class
  \IfFileExists{parskip.sty}{%
    \usepackage{parskip}
  }{% else
    \setlength{\parindent}{0pt}
    \setlength{\parskip}{6pt plus 2pt minus 1pt}}
}{% if KOMA class
  \KOMAoptions{parskip=half}}
\makeatother
\setlength{\emergencystretch}{3em} % prevent overfull lines
\providecommand{\tightlist}{%
  \setlength{\itemsep}{0pt}\setlength{\parskip}{0pt}}
\usepackage{amsmath,amssymb,mathtools,needspace}
\xeCJKDeclareCharClass{CJK}{"2032,"0400->"04FF,"2160->"216B,"2460->"2473,"25A0,"25C7,"25CB,"25CF}
\IfFontExistsTF{Arial Unicode MS}{\xeCJKsetup{AutoFallBack=true}\setCJKfallbackfamilyfont{\CJKrmdefault}{Arial Unicode MS}}{}
\setcounter{secnumdepth}{0}
\setlength{\emergencystretch}{2em}
\allowdisplaybreaks
\usepackage{bookmark}
\IfFileExists{xurl.sty}{\usepackage{xurl}}{} % add URL line breaks if available
\urlstyle{same}
\hypersetup{
  pdftitle={Gauss 点},
  colorlinks=true,
  linkcolor={Maroon},
  filecolor={Maroon},
  citecolor={Blue},
  urlcolor={Blue},
  pdfcreator={LaTeX via pandoc}}

\title{Gauss 点}
\author{}
\date{}

\begin{document}
\maketitle

\subsubsection{4.4.1 Gauss 点}\label{gauss-ux70b9}

Gauss 公式的求积节点称为 Gauss 点。

\textbf{定义 4.3}　如果求积公式（4.4.1）具有 \(2n+1\) 次代数精度，则称其节点 \(x_k\)（\(k=0,1,\cdots,n\)）是 Gauss 点。

下面从分析 Gauss 点的特性着手研究 Gauss 公式的构造问题。

\textbf{定理 4.4}　对于插值型求积公式（4.4.1），其节点 \(x_k\)（\(k=0,1,\cdots,n\)）是 Gauss 点的充分必要条件，是以这些点为零点的多项式 \(\displaystyle\omega(x)=\prod_{k=0}^{n}(x-x_k)\) 与任意次数不超过 \(n\) 的多项式 \(P(x)\) 均正交，即

\[
\int_a^b P(x)\omega(x)\,\mathrm{d}x=0.
\tag{4.4.2}
\]

\textbf{证明}　先证必要性。设 \(P(x)\) 是任意次数不超过 \(n\) 的多项式，则 \(P(x)\omega(x)\) 的次数不超过 \(2n+1\)。因此，如果 \(x_0,x_1,\cdots,x_n\) 是 Gauss 点，则求积公式（4.4.1）对于 \(P(x)\omega(x)\) 能准确成立，即有

\[
\int_a^b P(x)\omega(x)\,\mathrm{d}x
=\sum_{k=0}^{n}A_kP(x_k)\omega(x_k).
\]

但 \(\omega(x_k)=0\)（\(k=0,1,\cdots,n\)），故式（4.4.2）成立。

再证充分性。对于任意给定的次数不超过 \(2n+1\) 的多项式 \(f(x)\)，用 \(\omega(x)\) 除 \(f(x)\)，记商为 \(P(x)\)，余式为 \(Q(x)\)，\(P(x)\) 与 \(Q(x)\) 都是次数不超过 \(n\) 的多项式：

\[
f(x)=P(x)\omega(x)+Q(x).
\]

而利用式（4.4.2），得

\[
\int_a^b f(x)\,\mathrm{d}x=\int_a^b Q(x)\,\mathrm{d}x.
\tag{4.4.3}
\]

由于所给求积公式（4.4.1）是插值型的，它对于 \(Q(x)\) 能准确成立，即

\[
\int_a^b Q(x)\,\mathrm{d}x=\sum_{k=0}^{n}A_kQ(x_k).
\]

再注意到 \(\omega(x_k)=0\)，知 \(Q(x_k)=f(x_k)\)，从而有

\[
\int_a^b Q(x)\,\mathrm{d}x=\sum_{k=0}^{n}A_kf(x_k),
\]

于是由式（4.4.3）知

\[
\int_a^b f(x)\,\mathrm{d}x=\sum_{k=0}^{n}A_kf(x_k).
\]

（证明续下页。）

\href{../../source-images/pdf-108.jpeg}{原页图像}

\subsubsection{4.4.1 Gauss 点（续）}\label{gauss-ux70b9ux7eed}

可见求积公式（4.4.1）对于一切次数不超过 \(2n+1\) 的多项式均能准确成立，因此 \(x_k\)（\(k=0,1,\cdots,n\)）是 Gauss 点。证毕。

\begin{center}\rule{0.5\linewidth}{0.5pt}\end{center}

\subsection{相关原页校注（agent 补充）}\label{ux76f8ux5173ux539fux9875ux6821ux6ce8agent-ux8865ux5145}

以下校注来自本节覆盖的原页，可能也涉及同页相邻小节；原文照录与校正说明分开保留。

\subsubsection{原书 PDF 页 108 的校注}\label{ux539fux4e66-pdf-ux9875-108-ux7684ux6821ux6ce8}

\href{../pages/pdf-108.md}{完整原页转录} · \href{../../source-images/pdf-108.jpeg}{原页图像}

校注记录：CH04B-E002。

\begin{quote}
校注（agent 补充，CH04B-E002）：表 4.5 的 \(n=3\) 第二对节点在原扫描中印为 \(\pm0.339\,881\,0\)，本表原样保留；这是疑似原书数值排印错误。\(P_4(x)=(35x^4-30x^2+3)/8\) 的较小正零点为 \(\sqrt{(15-2\sqrt{30})/35}\approx0.339\,981\,0436\)，保留七位小数应为 \(0.339\,981\,0\)。原表权重 \(0.652\,145\,2\) 未改动。\href{../assets/ch04b/review-table-4-5.png}{表 4.5 原扫描局部}。
\end{quote}

\subsection{本节覆盖原页的校注记录}\label{ux672cux8282ux8986ux76d6ux539fux9875ux7684ux6821ux6ce8ux8bb0ux5f55}

下列记录按原页关联，含同页相邻小节的校注；计算时同时读取上文校注和完整原页。

\begin{itemize}
\tightlist
\item
  \texttt{CH04B-E002}：原书 PDF 页 108
\end{itemize}

\end{document}
